\documentclass[9pt,a4paper,twocolumn]{ctexart}

% --- Packages ---
\usepackage[top=1.2cm, bottom=1.2cm, left=1.2cm, right=1.2cm, columnsep=0.8cm]{geometry}
\usepackage{hyperref}
\usepackage{graphicx}
\usepackage{float}
\usepackage{longtable}
\usepackage{tabularx}
\usepackage{booktabs}
\usepackage{enumitem}
\usepackage{xcolor}
\usepackage{fancyhdr}
\usepackage{listings}
\usepackage{fontspec}
\usepackage{titlesec}
\usepackage{caption}

% --- Code block styling ---
\definecolor{codebg}{RGB}{245,245,245}
\definecolor{codeframe}{RGB}{220,220,220}
\definecolor{codekw}{RGB}{0,0,192}
\definecolor{codecomment}{RGB}{0,128,0}
\definecolor{codestring}{RGB}{163,21,21}
\definecolor{codenum}{RGB}{128,0,128}
\definecolor{codepunct}{RGB}{90,90,90}
\definecolor{badgegreen}{RGB}{34,139,34}
\definecolor{badgeblue}{RGB}{30,144,255}
\definecolor{badgeblack}{RGB}{50,50,50}

% --- Difficulty badges (rounded corners via tikz, pre-rendered in saveboxes) ---
\usepackage{tikz}
\newsavebox{\badgechubox}
\savebox{\badgechubox}{\tikz[baseline]{\node[fill=badgegreen,text=white,rounded corners=2.5pt,inner xsep=4pt,inner ysep=1.5pt,font=\sffamily\footnotesize\bfseries]{初级};}}
\newsavebox{\badgezhongbox}
\savebox{\badgezhongbox}{\tikz[baseline]{\node[fill=badgeblue,text=white,rounded corners=2.5pt,inner xsep=4pt,inner ysep=1.5pt,font=\sffamily\footnotesize\bfseries]{中级};}}
\newsavebox{\badgegaobox}
\savebox{\badgegaobox}{\tikz[baseline]{\node[fill=badgeblack,text=white,rounded corners=2.5pt,inner xsep=4pt,inner ysep=1.5pt,font=\sffamily\footnotesize\bfseries]{高级};}}
\newcommand{\lvchu}{\usebox{\badgechubox}}
\newcommand{\lvzhong}{\usebox{\badgezhongbox}}
\newcommand{\lvgao}{\usebox{\badgegaobox}}
\pdfstringdefDisableCommands{%
  \def\lvchu{[初级]}%
  \def\lvzhong{[中级]}%
  \def\lvgao{[高级]}%
}

\lstset{
  basicstyle=\ttfamily\scriptsize,
  keywordstyle=\color{codekw}\bfseries,
  commentstyle=\color{codecomment}\itshape,
  stringstyle=\color{codestring},
  backgroundcolor=\color{codebg},
  frame=single,
  rulecolor=\color{codeframe},
  breaklines=true,
  breakatwhitespace=false,
  breakindent=0pt,
  postbreak=\mbox{\textcolor{red}{$\hookrightarrow$}\space},
  numbers=left,
  numberstyle=\tiny\color{gray},
  columns=flexible,
  keepspaces=true,
  showstringspaces=false,
  tabsize=2,
  aboveskip=4pt,
  belowskip=4pt,
  extendedchars=true,
  literate={，}{{，}}1 {；}{{；}}1 {！}{{！}}1 {？}{{？}}1 {“}{{“}}1 {”}{{”}}1 {（}{{（}}1 {）}{{）}}1 {《}{{《}}1 {》}{{》}}1,
}

% --- Extra language definitions (no built-in listings support) ---
\lstdefinelanguage{json}{
  morekeywords={true,false,null},
  sensitive=true,
  morestring=[b]",
  comment=[l]{//},
  morecomment=[s]{/*}{*/},
}
\lstdefinelanguage{yaml}{
  morekeywords={true,false,null,yes,no,on,off},
  sensitive=false,
  morecomment=[l]{\#},
  morestring=[b]",
  morestring=[d]',
}
\lstdefinelanguage{http}{
  morekeywords={GET,POST,PUT,DELETE,PATCH,HEAD,OPTIONS,HTTP,Host,Accept,Authorization,Content-Type,Content-Length,User-Agent},
  sensitive=true,
  morecomment=[l]{\#},
  morestring=[b]",
}
\lstdefinelanguage{TypeScript}{
  morekeywords={abstract,any,as,async,await,break,case,catch,class,const,continue,debugger,declare,default,delete,do,else,enum,export,extends,false,finally,for,from,function,get,if,implements,import,in,instanceof,interface,keyof,let,module,namespace,never,new,null,number,of,package,private,protected,public,readonly,return,set,static,string,super,switch,symbol,this,throw,true,try,type,typeof,undefined,unknown,var,void,while,with,yield,Record,Array,Promise,AsyncIterable,ReadableStream,Date,Error,Object,JSON},
  sensitive=true,
  morecomment=[l]{//},
  morecomment=[s]{/*}{*/},
  morestring=[b]",
  morestring=[b]',
  morestring=[b]`{`},
}

% --- Compact spacing ---
\linespread{0.95}
\setlength{\parskip}{1pt plus 1pt minus 1pt}
\setlength{\parindent}{0pt}

% --- Table styling ---
\renewcommand{\arraystretch}{1.1}

% --- Heading styling (numbered) ---
\titleformat{\section}{\normalsize\bfseries}{\thesection\quad}{0em}{}
\titlespacing{\section}{0pt}{6pt}{2pt}
\titleformat{\subsection}{\small\bfseries}{\thesubsection\quad}{0em}{}
\titlespacing{\subsection}{0pt}{4pt}{1pt}
\titleformat{\subsubsection}{\small\bfseries}{\thesubsubsection\quad}{0em}{}
\titlespacing{\subsubsection}{0pt}{3pt}{1pt}

% --- Page style ---
\pagestyle{fancy}
\fancyhf{}
\fancyhead[L]{\small\emph{\leftmark}}
\fancyhead[R]{\small\thepage}
\renewcommand{\headrulewidth}{0.4pt}
\renewcommand{\sectionmark}[1]{}  % prevent sections from overwriting leftmark

% --- Hyperref setup ---
\hypersetup{
  colorlinks=true,
  linkcolor=blue,
  urlcolor=blue,
  citecolor=blue,
}

% --- Caption format ---
\DeclareCaptionFormat{myformat}{\small\bfseries #1#2#3}
\captionsetup{format=myformat}

\begin{document}

\title{LLM 运行机制：Token、上下文窗口与采样参数怎么影响输出}
\date{}
\maketitle
\markboth{Python 版 Agent 知识}{ Python 版 Agent 知识}

\section*{⭐️ Token 和上下文为什么决定成本与效果？}


当你在输入法里打“今天天气真”，它会自动建议“好”——大模型做的事情本质上一样。只不过它看的不是前面几个字，而是前面几千甚至几十万个字。每次只“补”一个 Token（文本碎片），然后把这个碎片加进上下文，再预测下一个，如此循环，直到生成完整回答。



这个过程叫做\textbf{自回归生成（Autoregressive Generation）}。



理解了自回归生成，后面所有概念都好办了：


\begin{itemize}[itemsep=1pt, topsep=2pt]
  \item \textbf{Token}：模型每一步“补”的文本碎片。
  \item \textbf{上下文窗口}：模型在“补”之前能看到多少文本。
  \item \textbf{Temperature / Top-p}：模型选哪个候选碎片的策略。
  \item \textbf{Max Tokens}：允许模型最多“补”多少步。
\end{itemize}


你可以把 Token 理解为“模型的阅读单位”。我们人类读中文是一个字一个字地看，读英文是一个词一个词地看。但模型既不按字、也不按词——它用一套自己的“拆字规则”（叫 Tokenizer）把文本切成大小不等的碎片，每个碎片就是一个 Token。



为什么不直接按字或按词切？因为模型需要在“词表大小”和“序列长度”之间取平衡：


\begin{itemize}[itemsep=1pt, topsep=2pt]
  \item 每个汉字都是一个 Token，词表小、但序列长（模型要“补”更多步）。
  \item 每个词都是一个 Token，序列短、但词表会爆炸（中文词组太多了）。
\end{itemize}


所以实际用的是折中方案——\textbf{子词切分算法}（如 BPE、Unigram），高频词保留为整体，低频词拆成更小片段。



你可以把 Token 想象成乐高积木。常用的“积木块”比较大（比如“你好”可能是一个 Token），不常用的词会被拆成更小的基础块拼起来。



Token 不是“一个字”或“一个词”的严格等价物：


\begin{itemize}[itemsep=1pt, topsep=2pt]
  \item 英文可能一个单词被拆成多个 Token。
  \item 中文可能一个词被拆成多个 Token，也可能多个字合并成一个 Token（取决于词频与词表）。
\end{itemize}


工程上通常用\textbf{经验估算}做容量规划，用\textbf{实际 API 返回的 usage}做精确计费与监控。



\textbf{经验估算（仅用于粗略规划）}：


\begin{itemize}[itemsep=1pt, topsep=2pt]
  \item 英文：1 Token 大约对应 3\textasciitilde\{\}4 个字符（与文本类型相关）。
  \item 中文：1 Token 常见在 1\textasciitilde\{\}2 个汉字上下波动（与混排比例强相关）。
\end{itemize}


DeepSeek 官方数据：1 个英文字符约消耗 0.3 Token，1 个中文字符约消耗 0.6 Token。换算过来，1 个 Token 约等于 3.3 个英文字符或 1.7 个中文字符，与上述经验值吻合。



成本趋势提示：Token 成本与 Tokenizer 版本强相关。早期模型（如 GPT-3.5）中文压缩率较低（约 1 字 1.5\textasciitilde\{\}2 Token）。GPT-4o 使用 o200k\_base Tokenizer（词表约 20 万），对中文压缩率有进一步提升；Qwen2.5 词表约 15 万，对中文常用词也有优化。实测数据因文本类型而异：新闻类约 1.5 字/Token，技术文档约 1.2 字/Token。



“趋近 1 字 1 Token”只适用于高频词汇，别拿它当成本估算基准。做预算前查一下当前模型版本的官方 Tokenizer 演示。



Token 划分直接影响模型理解能力。中文分词歧义和生僻字/低频专业术语的切分粒度，都会影响语义理解效果。



\textbf{Token 化过程示例}：


\begin{itemize}[itemsep=1pt, topsep=2pt]
  \item 原文：\texttt{你好，我是小 G。}
  \item 切分：\texttt{[你好]} \texttt{[，]} \texttt{[我是]} \texttt{[小 G]} \texttt{[。]}
  \item 统计：原文 9 字符 → Token 数 5 个 → 压缩比约 1.8 倍
\end{itemize}


\begin{figure}[H]
\centering
\includegraphics[width=0.85\textwidth]{/home/zhang/demos/latex_demos/Agent知识/HTML/LLM 运行机制：Token、上下文窗口与采样参数怎么影响输出 _ JavaGuide_files/llm-token-process.jpg}
\caption*{Token 化过程示例}
\end{figure}



注意：实际 Token 切分由模型供应商的 Tokenizer 实现，不同供应商对相同文本可能产生不同的 Token 序列。



OpenAI 官方网页端 Tokenizer 工具：\href{https://platform.openai.com/tokenizer}{OpenAI Tokenizer}



\textbf{特殊 Token}：除了文本内容对应的 Token，模型内部还会使用一些特殊标记，这些也会计入 Token 总数：


{\footnotesize
\begin{tabular}{|p{\dimexpr 0.441\columnwidth-2\tabcolsep\relax}|p{\dimexpr 0.322\columnwidth-2\tabcolsep\relax}|p{\dimexpr 0.237\columnwidth-2\tabcolsep\relax}|}
\hline
\textbf{特殊 Token} & \textbf{用途} & \textbf{示例} \\
\hline
BOS（Beginning of Sequence） & 标记序列开始 & \texttt{<s>} \\
\hline
EOS（End of Sequence） & 标记序列结束 & \texttt{</s>} \\
\hline
PAD（Padding） & 批处理时填充短序列 & \texttt{<pad>} \\
\hline
工具调用标记 & Function Calling 边界 & \texttt{<tool\_call/>} \\
\hline
\end{tabular}
}



这些特殊 Token 通常对用户不可见，但会占用上下文窗口。精确计数时建议使用官方 Tokenizer 工具而非手动估算。



\subsection*{多模态输入的 Token 开销}


GPT-4o、Claude 3.5、Gemini 等模型已支持图片输入。\textbf{图片不是“零成本”的}——它会被转换成一批 Token，同样占用上下文窗口。



粗略估算规则：


{\footnotesize
\begin{tabular}{|p{\dimexpr 0.128\columnwidth-2\tabcolsep\relax}|p{\dimexpr 0.423\columnwidth-2\tabcolsep\relax}|p{\dimexpr 0.449\columnwidth-2\tabcolsep\relax}|}
\hline
\textbf{模型} & \textbf{图片 Token 计算方式} & \textbf{一张 1024×1024 图片约等于} \\
\hline
GPT-4o & 按分辨率 + 细节模式 & 低细节 \textasciitilde\{\}85 tokens，高细节 \textasciitilde\{\}1105\textasciitilde\{\}765 tokens \\
\hline
Claude 3.5 & 固定 \textasciitilde\{\}5 tokens（缩略图）或 \textasciitilde\{\}85 tokens（全图） & 取决于图片模式 \\
\hline
Gemini & 按分辨率计算 & \textasciitilde\{\}258 tokens（标准） \\
\hline
\end{tabular}
}



工程启示：


\begin{itemize}[itemsep=1pt, topsep=2pt]
  \item 做多模态 RAG 时，要把图片 Token 也纳入预算。
  \item 批量处理图片时，注意首字延迟（TTFT）会显著增加。
  \item 如果只需要 OCR，考虑先用专门的 OCR 服务提取文字，再以纯文本形式送入模型。
\end{itemize}


\subsection*{上下文窗口的容量边界}


\textbf{上下文窗口}是 LLM 的“工作记忆”（Working Memory）。它决定了模型在任何时刻可以处理或“记住”的文本量（以 Token 为单位）。


\begin{itemize}[itemsep=1pt, topsep=2pt]
  \item 对话连续性：决定模型能进行多长的多轮对话而不遗忘早期细节。
  \item 单次处理能力：决定模型一次性能够处理的最大文档、代码库或数据样本。
\end{itemize}


“模型支持 128K/200K/1M”指的是一次调用里能放进模型的总 Token 上限。大多数模型的上下文窗口包含输入与输出的总和，但部分供应商（如 Google Gemini）对输入和输出分别设限，使用前请查阅具体 API 文档。



上下文窗口往往被隐形成本占用：



\begin{figure}[H]
\centering
\includegraphics[width=0.85\textwidth]{/home/zhang/demos/latex_demos/Agent知识/HTML/LLM 运行机制：Token、上下文窗口与采样参数怎么影响输出 _ JavaGuide_files/llm-context-window.jpg}
\caption*{上下文窗口（Context Window）= LLM 的「工作记忆」}
\end{figure}


\begin{itemize}[itemsep=1pt, topsep=2pt]
  \item System Prompt：调节模型行为的系统指令（对用户隐藏，但占用窗口）。
  \item User Prompt：业务数据与指令。
  \item 多轮对话历史：过往的消息记录。
  \item RAG 检索片段：从外部知识库检索到的补充信息。
  \item 工具调用 Schema：函数定义与参数结构。
  \item 格式开销：特殊字符、换行符、Markdown 标记等。
  \item 模型生成的输出 Token：\textbf{输出也占用上下文窗口}。
\end{itemize}


因此，你真正能塞进 Prompt 的“有效业务内容”往往远小于标称上限。



注意：上下文窗口（Context Window）≠ 最大生成长度。许多模型支持 128K 甚至 1M 输入，但单次输出上限因 API 而异。OpenAI Chat Completions API 使用 \texttt{max\_tokens} 参数（GPT-4o 最大 16K 输出），部分新模型支持 \texttt{max\_completion\_tokens}（如 o1 系列），DeepSeek V3 最大输出 8K。使用前需查阅具体模型的 API 文档。



思维链模式的多轮对话处理：思维链模型（如 DeepSeek-R1）的 \texttt{reasoning\_content}（思考过程）通常不会被自动包含在下一轮对话的上下文中，只有 \texttt{content}（最终回答）会参与后续对话。



这意味着：


\begin{itemize}[itemsep=1pt, topsep=2pt]
  \item 无需为思考过程额外占用上下文窗口。
  \item 如果后续对话需要参考之前的推理过程，需要手动将 \texttt{reasoning\_content} 拼接到消息历史中。
  \item 部分供应商的 SDK 会自动处理这一差异，建议查阅具体文档确认。
\end{itemize}


\subsection*{长上下文背后的计算约束}


上下文窗口并非越大越好，它受限于 Transformer 架构的\textbf{自注意力机制（Self-Attention）}：


\begin{itemize}[itemsep=1pt, topsep=2pt]
  \item 计算成本平方级增长：计算需求与序列长度呈平方级关系（O(N²)）。输入 Token 翻倍，处理能力需求可能变为 4 倍。
  \item 推理延迟增加：上下文变长后，模型生成每个新 Token 时需要关注的历史 Token 变多，首字延迟 TTFT 会显著增加。
  \item 安全风险增加：更长的上下文意味着更大的攻击面。
\end{itemize}


工程优化手段：FlashAttention、GQA/MQA、Sliding Window Attention、Ring Attention 等技术已显著降低长上下文的计算和显存开销。但 O(N²) 的理论复杂度仍是上限扩展的根本瓶颈。



\subsection*{上下文溢出的真实表现}


当上下文接近上限或内容过长时，常见现象包括：


\begin{itemize}[itemsep=1pt, topsep=2pt]
  \item 模型忽略早期约束：System Prompt 里要求“必须输出 JSON”，但因距离生成点太远，注意力不足导致被忽略。
  \item “中间丢失”现象：即使在 1M 窗口模型中，模型对开头和结尾的信息最敏感，对中间部分的信息召回率显著下降。
  \item 回答漂移：前半段还围绕问题，后半段开始总结/扩写/跑题。
  \item RAG 失效：检索文档过多，关键信息被稀释；或被截断导致证据链断裂。
  \item 成本与延迟激增：1M 上下文会导致 TTFT 显著增加，且 Token 成本呈线性增长。
\end{itemize}


\subsection*{输入 Token 与输出 Token 的计费差异}


大多数供应商对输入 Token 和输出 Token 采用不同的计费标准，通常输出价格是输入的 \textbf{2\textasciitilde\{\}4 倍}：


{\footnotesize
\begin{tabular}{|p{\dimexpr 0.309\columnwidth-2\tabcolsep\relax}|p{\dimexpr 0.291\columnwidth-2\tabcolsep\relax}|p{\dimexpr 0.291\columnwidth-2\tabcolsep\relax}|p{\dimexpr 0.109\columnwidth-2\tabcolsep\relax}|}
\hline
\textbf{模型} & \textbf{输入价格（/1M Tokens）} & \textbf{输出价格（/1M Tokens）} & \textbf{输出/输入比} \\
\hline
GPT-4o & \$2.50 & \$10.00 & 4x \\
\hline
Claude 3.5 Sonnet & \$3.00 & \$15.00 & 5x \\
\hline
DeepSeek V3 & ¥0.5 & ¥2.0 & 4x \\
\hline
DeepSeek-R1 & ¥4.0 & ¥16.0 & 4x \\
\hline
\end{tabular}
}



工程启示：


\begin{itemize}[itemsep=1pt, topsep=2pt]
  \item 长 Prompt + 短输出 = 更经济的调用方式。
  \item RAG 场景要控制检索片段数量，避免输入 Token 激增。
  \item 思维链模型的 reasoning tokens 通常按输出价格计费，成本更高。
\end{itemize}


\subsection*{Prompt Caching 的省钱逻辑}


当请求中存在大量重复的固定前缀（如 System Prompt、长 RAG Context），可以用 \textbf{Prompt Caching} 显著降低成本。



原理：供应商会缓存请求中“可复用的前缀部分”。下次请求如果前缀相同，这部分就不重新计费，只收“缓存读取”的费用（通常是正常价格的 10\%\textasciitilde\{\}50\%）。



典型适用场景：


\begin{itemize}[itemsep=1pt, topsep=2pt]
  \item 多轮对话（System Prompt + 历史 Message 不变）。
  \item RAG 应用（检索片段重复率高）。
  \item 批量评估（同一份 System Prompt，不同的简历/文章）。
\end{itemize}


各供应商支持情况：


{\footnotesize
\begin{tabular}{|p{\dimexpr 0.220\columnwidth-2\tabcolsep\relax}|p{\dimexpr 0.366\columnwidth-2\tabcolsep\relax}|p{\dimexpr 0.195\columnwidth-2\tabcolsep\relax}|p{\dimexpr 0.220\columnwidth-2\tabcolsep\relax}|}
\hline
\textbf{供应商} & \textbf{功能名称} & \textbf{缓存时长} & \textbf{缓存命中折扣} \\
\hline
OpenAI & Prompt Caching & 5\textasciitilde\{\}10 分钟 & 输入价格约 50\\% \\
\hline
Anthropic & Prompt Caching & 5 分钟 & 输入价格约 10\\% \\
\hline
DeepSeek & Context Caching & 10\textasciitilde\{\}30 分钟 & 输入价格约 25\\% \\
\hline
\end{tabular}
}



工程建议：


\begin{enumerate}[itemsep=1pt, topsep=2pt]
  \item 把不变的内容放前面（System Prompt、工具定义、RAG Context），把变化的内容放后面（User Prompt）。
  \item 监控 \texttt{cache\_read\_tokens} 和 \texttt{cache\_creation\_tokens} 指标，验证缓存命中率。
  \item 批量任务尽量在缓存时间窗口内完成。
\end{enumerate}


\subsection*{一次调用的 Token 预算公式}


把“上下文窗口”当成一个固定容量的桶，下图展示了一个典型调用的 Token 预算分配：



此分配仅为示意，实际比例需根据业务场景动态调整。



最实用的预算方式是：



\textbf{window ≥ input\_tokens + max\_output\_tokens}



对于思维链模型，公式应调整为：



\textbf{window ≥ input\_tokens + reasoning\_tokens + max\_output\_tokens}



其中 \texttt{reasoning\_tokens}（思考链 Token 数）难以精确预估，建议按 \texttt{max\_output\_tokens} 的 2\textasciitilde\{\}3 倍预留。



其中 \texttt{input\_tokens} 至少包含：


\begin{itemize}[itemsep=1pt, topsep=2pt]
  \item system prompt（含 schema / 工具定义）
  \item user prompt（含变量替换后的实际文本）
  \item 历史消息（多轮对话时）
  \item RAG context（如果拼进来了）
\end{itemize}


工程上建议反过来做预算（因为输出经常更可控）：


\begin{enumerate}[itemsep=1pt, topsep=2pt]
  \item 先定 \texttt{max\_output\_tokens}（结构化输出通常不需要很长）。
  \item 再为输入预留安全边际（例如再留 10\%\textasciitilde\{\}20\% 给供应商额外开销）。
  \item 超预算时，用可解释的策略“减输入”而不是“赌模型会自我约束”：
\end{enumerate}
\begin{itemize}[itemsep=1pt, topsep=2pt]
  \item 优先减少 RAG 的 Top-K 或做片段去重。
  \item 对长字段做摘要/截断（如简历、长回答）。
  \item 多段任务拆成多次调用（分批评估、两阶段生成）。
\end{itemize}


\section*{⭐️ 采样参数如何影响输出稳定性？}


\subsection*{从 logits 到概率采样}


模型每一步会给词表中\textbf{每个}候选 Token 打一个分数（内部叫 \textbf{logits}），分数越高说明模型越觉得这个词应该出现在这里。



举个例子，假设模型正在补全“今天天气真\_\_”，它可能给出这样的分数：


{\footnotesize
\begin{tabular}{|p{\dimexpr 0.421\columnwidth-2\tabcolsep\relax}|p{\dimexpr 0.579\columnwidth-2\tabcolsep\relax}|}
\hline
\textbf{候选 Token} & \textbf{原始分数（logit）} \\
\hline
好 & 5.0 \\
\hline
不错 & 3.2 \\
\hline
棒 & 2.1 \\
\hline
糟糕 & 0.5 \\
\hline
紫色 & -8.0 \\
\hline
\end{tabular}
}



但原始分数不是概率——需要经过一次数学变换（\textbf{softmax}）才能变成每个候选被选中的概率。变换后大致是：


{\footnotesize
\begin{tabular}{|p{\dimexpr 0.667\columnwidth-2\tabcolsep\relax}|p{\dimexpr 0.333\columnwidth-2\tabcolsep\relax}|}
\hline
\textbf{候选 Token} & \textbf{概率} \\
\hline
好 & 62\\% \\
\hline
不错 & 20\\% \\
\hline
棒 & 10\\% \\
\hline
糟糕 & 5\\% \\
\hline
紫色 & ≈ 0\\% \\
\hline
\end{tabular}
}



最后，模型按这个概率分布“抽签”（采样），决定输出哪个 Token。



解码参数（Temperature、Top-p、Top-k 等）就是在这个“打分 → 概率 → 抽签”的过程中施加控制：


\begin{itemize}[itemsep=1pt, topsep=2pt]
  \item Temperature：调整概率分布的“形状”，让高分选项更突出，或者让各选项更均匀。
  \item Top-p / Top-k：直接砍掉不靠谱的候选项，缩小“抽签池”。
  \item Penalty 系列：对已经出现过的词降分，防止“复读机”。
\end{itemize}


\subsection*{Temperature 的“冒险程度”}


\begin{figure}[H]
\centering
\includegraphics[width=0.85\textwidth]{/home/zhang/demos/latex_demos/Agent知识/HTML/LLM 运行机制：Token、上下文窗口与采样参数怎么影响输出 _ JavaGuide_files/llm-temperature-params.jpg}
\caption*{Temperature 参数：控制模型输出的随机性}
\end{figure}



Temperature 的工作原理很简单：在 softmax 之前，先把所有分数\textbf{除以}温度值 T。



\textbf{p(t) = softmax(z\_t / T)}


\begin{itemize}[itemsep=1pt, topsep=2pt]
  \item T ≈ 1：保持原始分布。
  \item T < 1：分布更尖锐，更倾向选择高概率 Token（更“稳”）
  \item T > 1：分布更平坦，低概率 Token 更容易被采样到（更“野”）
\end{itemize}


还是用“今天天气真\_\_”的例子：


\begin{itemize}[itemsep=1pt, topsep=2pt]
  \item T = 0.2（低温）：分数差距被放大（都除以 0.2，等于乘以 5），原本就领先的“好”概率飙升到 \textasciitilde\{\}98\%，几乎每次都选它。
  \item T = 1.0（默认温度）：保持原始分布不变，“好”62\%、“不错”20\%...按正常概率采样。
  \item T = 1.5（高温）：分数差距被缩小（都除以 1.5），“好”概率降到 \textasciitilde\{\}35\%，“棒”、“不错”甚至“糟糕”都有更大机会被选中。
\end{itemize}


温度越低，输出越确定；温度越高，输出越随机。



工程建议（经验值，非硬规则）：


{\footnotesize
\begin{tabular}{|p{\dimexpr 0.319\columnwidth-2\tabcolsep\relax}|p{\dimexpr 0.213\columnwidth-2\tabcolsep\relax}|p{\dimexpr 0.468\columnwidth-2\tabcolsep\relax}|}
\hline
\textbf{场景} & \textbf{推荐温度} & \textbf{说明} \\
\hline
结构化提取 / JSON 输出 & 0 \textasciitilde\{\} 0.3 & 配合严格 schema + 解析失败重试策略 \\
\hline
评估 / 分析 / 代码评审 & 0.4 \textasciitilde\{\} 0.8 & 平衡确定性与表达多样性 \\
\hline
创作类内容（文案、头脑风暴） & 0.8 \textasciitilde\{\} 1.2+ & 增加多样性，但要承担格式一致性风险 \\
\hline
\end{tabular}
}



追求确定性？若需单元测试幂等或结果复现，仅设 \texttt{Temperature=0} 不够（GPU 浮点误差仍可能导致非确定性）。建议同时配置 \textbf{\texttt{seed} 参数}（如 OpenAI/DeepSeek 支持）。



即使配置 \texttt{seed}，以下情况仍可能导致结果不一致：


\begin{itemize}[itemsep=1pt, topsep=2pt]
  \item 模型版本更新（底层权重变化）。
  \item 跨区域调用（不同集群可能部署不同版本）。
  \item Top-p 采样（即使 T=0，若 Top-p<1 仍有随机性）。
\end{itemize}


建议在 CI/CD 中仅将 LLM 调用用于冒烟测试，核心逻辑仍依赖 Mock。



\subsection*{Top-p 与 Top-k 的“抽签池”}


Temperature 调整的是概率分布的形状，但不管怎么调，词表里所有 Token 理论上都有被选中的可能。Top-p 和 Top-k 则更直接——把不靠谱的候选直接踢出抽签池。



还是用“今天天气真\_\_”的例子：


{\footnotesize
\begin{tabular}{|p{\dimexpr 0.500\columnwidth-2\tabcolsep\relax}|p{\dimexpr 0.188\columnwidth-2\tabcolsep\relax}|p{\dimexpr 0.312\columnwidth-2\tabcolsep\relax}|}
\hline
\textbf{候选 Token} & \textbf{概率} & \textbf{累计概率} \\
\hline
好 & 62\\% & 62\\% \\
\hline
不错 & 20\\% & 82\\% \\
\hline
棒 & 10\\% & 92\\% \\
\hline
糟糕 & 5\\% & 97\\% \\
\hline
紫色 & ≈0\\% & ≈100\\% \\
\hline
\end{tabular}
}


\begin{itemize}[itemsep=1pt, topsep=2pt]
  \item Top-k = 3：只保留概率最高的 3 个候选（好、不错、棒），在这 3 个里重新分配概率后采样。“糟糕”和“紫色”直接出局。
  \item Top-p = 0.9：从高到低累加概率，保留累计刚好达到 90\% 的最小集合。这里“好 + 不错 + 棒 = 92\% ≥ 90\%”，所以保留这 3 个。如果某个场景下头部更集中（比如第一名就占了 95\%），Top-p 会自动只保留 1 个——比 Top-k 更灵活的地方就在这。
\end{itemize}


两者的区别：Top-k 固定保留 k 个，不管概率分布长什么样；Top-p 根据概率自适应调整候选数量。实践中 \textbf{Top-p 更常用}，因为它能自动适应不同的概率分布。



常见组合：


{\footnotesize
\begin{tabular}{|p{\dimexpr 0.372\columnwidth-2\tabcolsep\relax}|p{\dimexpr 0.372\columnwidth-2\tabcolsep\relax}|p{\dimexpr 0.256\columnwidth-2\tabcolsep\relax}|}
\hline
\textbf{组合} & \textbf{效果} & \textbf{适用场景} \\
\hline
T=0（贪婪解码） & 永远选最高分，完全确定 & 结构化输出、可复现场景 \\
\hline
低温 + Top-p=0.9 & 相对稳定，但允许措辞上有些变化 & 分析报告、摘要 \\
\hline
中高温 + Top-p=0.95 & 多样性较高，但排除了极端离谱选项 & 创意写作、对话 \\
\hline
\end{tabular}
}



注意：贪婪解码虽然最稳定，但可能更容易陷入重复循环。



\subsection*{停止条件与截断风险}


工程上需要意识到两点：


\begin{itemize}[itemsep=1pt, topsep=2pt]
  \item \textbf{Max Tokens 是硬上限}：到上限会被强制截断，模型正写到一半也会被“掐断”。常见后果：JSON 缺右括号、列表缺最后几项、句子写了一半。
  \item \textbf{Stop Sequences（停止词）是软切断}：可以指定一些字符串（如 \texttt{"\textbackslash\{\}n\textbackslash\{\}n"} 或 ```\texttt{ "}`\texttt{" }```），模型生成到这些内容时会自动停止。但如果 stop 设计不当，可能提前截断关键字段。
\end{itemize}


结构化输出场景要把“截断风险”当成一类失败路径来设计缓解策略。



思维链模式的 Token 计算差异：对于支持思维链的模型（如 DeepSeek-R1），\texttt{max\_tokens} 通常包含思考过程 + 最终回答两部分。例如设置 \texttt{max\_tokens=8192}，模型可能在思考链上消耗 5000 tokens，最终回答只剩 3192 tokens 的预算。



不同供应商的默认值和上限差异较大：DeepSeek-R1 默认 32K、最大 64K；OpenAI o1 系列的输出上限也高于普通模型。使用前务必查阅具体模型的 API 文档。



\subsection*{Penalty 与复读问题}


可能遇到过模型反复输出同一句话，或者在长回答里不断重复相同观点。Penalty 参数用来缓解这类问题，它们在解码时\textbf{降低已出现 Token 的分数}：


{\footnotesize
\begin{tabular}{|p{\dimexpr 0.346\columnwidth-2\tabcolsep\relax}|p{\dimexpr 0.404\columnwidth-2\tabcolsep\relax}|p{\dimexpr 0.250\columnwidth-2\tabcolsep\relax}|}
\hline
\textbf{参数} & \textbf{作用} & \textbf{通俗理解} \\
\hline
Repetition Penalty & 降低所有已出现 Token 的概率 & “说过的词，再说就扣分” \\
\hline
Presence Penalty & 只要 Token 出现过就扣分（不看次数） & “鼓励聊新话题” \\
\hline
Frequency Penalty & Token 出现次数越多扣分越重 & “同一个词说了三遍？重罚” \\
\hline
\end{tabular}
}



工程陷阱：


\begin{itemize}[itemsep=1pt, topsep=2pt]
  \item 结构化输出别乱加 Penalty：JSON 里字段名（如 \texttt{"name"}、\texttt{"score"}）需要反复出现，加了 Repetition Penalty 可能把必须出现的字段名也“惩罚掉”，导致输出残缺。
  \item RAG 问答别加 Presence Penalty：它会鼓励模型“说点新东西”，反而降低对检索内容的忠实度，增加幻觉风险。
\end{itemize}


保守建议：如果不确定这些参数的精确语义（不同供应商定义可能不同），建议保持默认值。用低温 + 更强 Prompt 约束 + 更短输出来获得稳定性，比调 Penalty 更可控。



\subsection*{思维链模式的参数限制}


部分模型（如 DeepSeek-R1、OpenAI o1）支持“思维链模式”，在生成最终回答前会先输出一段内部推理过程。这类模型有特殊的参数约束：



不支持的采样参数：思维链模式下，以下参数通常被忽略：


\begin{itemize}[itemsep=1pt, topsep=2pt]
  \item \texttt{temperature}、\texttt{top\_p}：采样控制参数。
  \item \texttt{presence\_penalty}、\texttt{frequency\_penalty}：惩罚参数。
\end{itemize}


原因：思维链模式的设计目标是让模型“自由思考”，采用模型内部固定的采样策略，用户传入的采样参数会被忽略。



工程建议：


\begin{itemize}[itemsep=1pt, topsep=2pt]
  \item 调用思维链模型时，不要依赖上述参数控制输出风格。
  \item 若需要更稳定的输出格式，应通过 Prompt 约束而非采样参数。
  \item 关注模型返回的 \texttt{reasoning\_content} 字段（思考过程）与 \texttt{content} 字段（最终回答）的区别。
\end{itemize}


\subsection*{流式输出与首字延迟}


默认情况下，API 会等模型生成完所有内容后一次性返回。流式输出则是边生成边返回——模型每生成一个（或几个）Token，就立刻推送给客户端，用户更早看到内容开始出现。



核心价值：改善用户体验，降低首字延迟（TTFT，Time-To-First-Token）。



常见误解澄清：


\begin{itemize}[itemsep=1pt, topsep=2pt]
  \item 流式输出更快——总耗时（E2E latency）不一定下降，模型生成的总 Token 量相同。
  \item 流式输出更省钱——Token 计费不变，仍然受限流/配额影响。
  \item 如果需要结构化输出（如 JSON），流式场景要考虑“半成品 JSON”在前端/网关层的处理。
\end{itemize}


\subsection*{Logprobs 与置信度排查}


部分 API（如 OpenAI）支持返回每个生成 Token 的\textbf{对数概率}（logprobs），可以理解为模型对该 Token 的“确信程度”。logprob 越接近 0，模型越确信；值越小（如 -5.0），说明模型越“犹豫”。



工程应用场景：


\begin{itemize}[itemsep=1pt, topsep=2pt]
  \item \textbf{置信度评估}：提取“金额: 1000”时，若对应 Token 的 logprob 很低，说明模型不太确定，可能需要人工复核。
  \item \textbf{异常检测}：监控生产环境中模型输出的平均 logprob，若突然下降可能提示 Prompt 漂移或输入数据异常。
  \item \textbf{多候选对比}：获取 Top-N 候选 Token 及其概率，用于纠错或二次排序。
\end{itemize}


注意事项：logprobs 会增加响应体积，且并非所有供应商都支持。使用前请查阅 API 文档。



\subsection*{采样参数配置建议}

{\footnotesize
\begin{tabular}{|p{\dimexpr 0.214\columnwidth-2\tabcolsep\relax}|p{\dimexpr 0.196\columnwidth-2\tabcolsep\relax}|p{\dimexpr 0.089\columnwidth-2\tabcolsep\relax}|p{\dimexpr 0.125\columnwidth-2\tabcolsep\relax}|p{\dimexpr 0.375\columnwidth-2\tabcolsep\relax}|}
\hline
\textbf{场景} & \textbf{Temperature} & \textbf{Top-p} & \textbf{Penalty} & \textbf{其他建议} \\
\hline
JSON / 结构化输出 & 0 \textasciitilde\{\} 0.3 & 1.0 & 保持默认 & 配合 Strict Mode + 重试策略 \\
\hline
代码评审 / 技术分析 & 0.4 \textasciitilde\{\} 0.7 & 0.9 & 保持默认 & 结合 CoT Prompt \\
\hline
多轮对话 & 0.6 \textasciitilde\{\} 0.8 & 0.9 & 适度开启 & 控制历史消息长度 \\
\hline
创意写作 / 头脑风暴 & 0.8 \textasciitilde\{\} 1.2 & 0.95 & 按需开启 & 接受输出多样性，做好后处理 \\
\hline
思维链模型 & —（不支持） & — & — & 通过 Prompt 控制，非采样参数 \\
\hline
\end{tabular}
}



\section*{总结}


回顾这篇扫盲内容，核心就是处理好三个维度的工程权衡：


\begin{enumerate}[itemsep=1pt, topsep=2pt]
  \item \textbf{Token 是成本与性能的物理标尺}：它不仅决定计费账单和推理延迟，更决定模型对文本的理解粒度。做容量规划时，必须按 Token 算账，而不是按字数算账。
  \item \textbf{上下文窗口是极其稀缺的资源}：哪怕模型宣称支持 1M 上下文，也不意味着可以毫无节制地堆砌数据。为 Prompt、RAG 检索片段、历史对话和输出预留做好严格的 Token 预算分配，是走向生产环境的必修课。
  \item \textbf{采样参数是业务场景的调音台}：如果追求稳定的 JSON 输出，就果断压低 Temperature 并配合严格的 Schema；如果需要创意与头脑风暴，再适度放开 Temperature 和 Top-p。不要迷信默认参数，要根据业务的容错率来定制。
\end{enumerate}


打好这层参数与原理的地基，再去看 Agent 编排、RAG 检索或是 MCP 工具调用，你会发现那些高阶架构的本质，无非是在更好地调度这些底层 Token，更精准地管理这个上下文窗口。


\end{document}
